\chapter{The mixed sectors}
\label{chap:mixed}

By Lemma~\ref{lem:sector-swap} it suffices to treat the sector $(-,+)$: odd
under $S_1$ (Dirichlet on the edges $a=\pm\pi/2$) and even under $S_2$ (Neumann
on the edges $b=\pm\pi/2$). We write $D$ for the Dirichlet and $N$ for the
Neumann edges of $\Cell$. Throughout this chapter eigenfunctions are at
real $\lambda\ne0$ in the sector $(-,+)$. We use the continuous representative
of $\phi=\psi/\cos a$ from Proposition~\ref{prop:mixed-quotient-domain}
for its edge and corner values.

\emph{Local notation:} $a=(x+y)/2$, $b=(x-y)/2$, $h=\cos a$, $k=\cos b$,
$u=\sin a$, $v=\sin b$, $s=\Psi=2hk$, $p=h+k$, $R=u^2+v^2$, $\mu=h^2$.
Here $dx\,dy=2\,da\,db$ and the physical gradient in the orthonormal
diagonal directions is $\nabla_0/\sqrt2$, where $\nabla_0=(\partial_a,\partial_b)$.
The edge subscripts $n,t$ denote physical unit normal and tangent
derivatives; $t$ as an integration variable is directed flow time.

\begin{remark}[Digest note: notation]
Within this sector the manuscript uses three overlapping naming schemes for
the same objects: $(G,V,g,j)$ in the completed square, $(\mathsf G,U,D,E)$ in
the comparison lemma (with $U=-V$), and $(G,V,d,c,S,R)$ in the comparison
certificates (with $S=p$ and $R=2-h^2-k^2=u^2+v^2$). The letter $w$ means
$1-x^2$ on the edges but the comparison function in the certificates. This
should be unified in the paper.
\end{remark}

\section{Virial identities}

\begin{definition}[Mixed field]
\label{def:mixed-field}
\uses{def:cells, lem:cell-square, lem:diagonal-calculus, def:adapted-field}
\lean{NoEmbeddedEigenvalues.mixedFieldX,
 NoEmbeddedEigenvalues.mixedFieldY,
 NoEmbeddedEigenvalues.mixedFieldA,
 NoEmbeddedEigenvalues.mixedFieldB,
 NoEmbeddedEigenvalues.mixedFieldCoeff}
\leanok
On $\Cell$ let
$W=\tan(x/2)\partial_x+\tan(y/2)\partial_y=\frac{u}{p}\partial_a+\frac{v}{p}\partial_b$
and $f=R/p^2$.
\end{definition}

\begin{lemma}[Mixed field algebra]
\label{lem:mixed-field-algebra}
\uses{def:mixed-field}
\lean{NoEmbeddedEigenvalues.mixed_field_algebra, NoEmbeddedEigenvalues.mixedFieldMatrixOffDiagonal}
\leanok
$W\Psi=\Psi-2$, $W$ is adapted with coefficient $f$, $\Div W=f+1$, and in the
orthonormal directions parallel to the $a$, $b$ axes
\[
  M:=2\sym DW-fI=\begin{pmatrix}1&2uv/p^2\\2uv/p^2&1\end{pmatrix}.
\]
The eigenvalues of $2\sym DW$ are $\sec^2((a+b)/2)$ and $\sec^2((a-b)/2)$.
\end{lemma}

\begin{proof}
\uses{lem:dipole-identities}
Direct differentiation. In $U=\sin(x/2)$, $V=\sin(y/2)$ the field is
$(U\partial_U+V\partial_V)/2$.
\end{proof}

\begin{proposition}[Paired flux identity for the normal field]
\label{prop:mixed-paired-flux}
\uses{def:normal-field, def:canonical-profile, prop:mixed-quotient-domain}
\lean{NoEmbeddedEigenvalues.mixed_paired_flux}
\leanok
Write $f_0,c_0$ for the coefficients in Definition~\ref{def:normal-field}.
For a mixed eigenfunction and its quotient,
\[
  \int_\Cell\big[f_0|\omega|^2+(2-c_0)|\psi_n|^2+c_0|\psi_t|^2\big]\,dx\,dy
  =2\lim_{T\to\infty}\Big\{\sum_D\int_{-T}^T|\psi_n|^2\,dt
  -\sum_N\int_{-T}^T\big[|\psi_t|^2+\operatorname{Re}(\Omega\bar g)\big]\,dt\Big\},
\]
where on a Neumann edge $g=\psi|_N$ and $\Omega$ is its canonical profile. The
bulk integral and the paired limit are finite although individual edge
energies may diverge (manuscript \texttt{(mf:paired)}).
\end{proposition}

\begin{proof}
\uses{lem:ee-truncated-virial, prop:saddle-first-gain, cor:mixed-saddle-jet,
  lem:saddle-normal-form, lem:pairing-limit}
The bulk is integrable since $f_0\le C/d^2$ and $\omega/d\in L^2$. Apply
the truncated virial along good levels. In inward orthonormal saddle
coordinates $(\xi,\eta)$ the model jet $\ell=A\xi$ has
\[
 (|\ell_n|^2-|\ell_t|^2)\,dt
  = |A|^2\frac{\eta^2-\xi^2}{\xi^2+\eta^2}\frac{d\xi}{\xi}
 \quad\text{on }\xi\eta=\epsilon.
\]
Its integral on a symmetrically truncated hyperbola vanishes under
$\xi\leftrightarrow\eta$, cancelling the two adjacent logarithms.
The weighted jet estimate controls the remainder. For the canonical
pairing, apply Lemma~\ref{lem:pairing-limit} to both character components
on a common good sequence. The cross pairings cancel: shifting by one
quarter multiplies the product of a $+i$ component and the conjugate of a
$-i$ component by $-1$, both on a regular orbit and in the sum on the arcs.
\end{proof}

\begin{remark}[Digest note: two parts of the mixed flux limit]
Lemma~\ref{lem:pairing-limit} now supports the $\pm i$ characters on a
prescribed good sequence carrying both transports. The mixed application
still has to choose that sequence and formalize the recombination.
Separately, the even--even gradient trace theorem cannot handle the
signed corner cancellation: mixed linear jets can make the individual
edge energies divergent.
\end{remark}

\begin{proposition}[Mixed virial identity]
\label{prop:mixed-virial}
\uses{def:mixed-field, prop:mixed-quotient-domain}
\lean{NoEmbeddedEigenvalues.mixed_virial}
\leanok
For a mixed eigenfunction and its quotient, with the same centered flow-time
cutoff on all four edges,
\[
  \int_\Cell \big(f|\omega|^2+\ip{M\nabla\psi}{\nabla\psi}\big)\,dx\,dy
  =2\lim_{T\to\infty}\Big\{\sum_D\int_{-T}^T|\psi_n|^2\,dt
  -\sum_N\int_{-T}^T\big(|\partial_\tau\psi|^2+\operatorname{Re}(\omega|_N\,\bar\psi)\big)dt\Big\},
\]
and the bulk integral converges absolutely (manuscript
\texttt{(prop:r5-new-virial)}).
\end{proposition}

\begin{proof}
\uses{prop:mixed-paired-flux, lem:mixed-field-algebra, lem:mixed-boundary-square-traces,
  cor:mixed-saddle-jet, prop:saddle-first-gain}
$W-W_0$, with $W_0$ the normal field, is tangent to the orbits, and the
tangential virial on $\{s>\epsilon\}$ gives
$J_W-J_{W_0}=2\operatorname{Re}\int_{\partial}((W-W_0)\psi)\,\partial_n\bar\psi\,d\ell$.
This current vanishes on regular edges by the Dirichlet or Neumann condition,
and its corner contribution
$4|A|^2\frac{UV(U-V)}{(U+V)(U^2+V^2)}\frac{dU}U$ integrates to zero by
$U\leftrightarrow V$ symmetry.
\end{proof}

\section{The boundary square}

\begin{lemma}[Mixed boundary square, smooth case]
\label{lem:mixed-boundary-square}
\lean{NoEmbeddedEigenvalues.mixed_boundary_square,
 NoEmbeddedEigenvalues.mixedBoundaryTransport,
 NoEmbeddedEigenvalues.mixedBoundarySquare}
\leanok
With $w=1-x^2$, $T_\lambda F=-i(wF'-xF)-\lambda F$ and $G=T_\lambda F$, for
smooth $F$ on $[-1,1]$ and real $\lambda$,
\[
  \int_{-1}^1w|G'|^2\,dx-2\operatorname{Re}\int_{-1}^1T_0F\,\bar G\,dx
  =\int_{-1}^1w\,|wF''-3xF'-i\lambda F'|^2\,dx
\]
(manuscript \texttt{(eq:mixed-boundary-square)}). No boundary conditions on
$F$ are needed.
\end{lemma}

\begin{proof}
With $Sv=-i(wv'-3xv)-\lambda v$ one has $G'=S(F')+iF$. Expanding, the
difference of the two sides is
$\int[(1-3x^2)|F|^2-2w^2|F'|^2-2w^2\operatorname{Re}(F''\bar F)
+10xw\operatorname{Re}(F'\bar F)]$, the $\lambda$ terms cancelling before
integration; integrating the third term by parts leaves
$\int[(1-3x^2)|F|^2+xw(|F|^2)']=0$. All boundary currents carry a factor $w$.
\end{proof}

\begin{lemma}[Mixed boundary square on the weighted domain]
\label{lem:mixed-boundary-square-weighted}
\uses{lem:mixed-boundary-square, def:h1-interval}
\lean{NoEmbeddedEigenvalues.mixed_boundary_square_weighted}
\leanok
Let $\lambda\in\R$, $w=1-x^2$, and let $F,G$ be locally absolutely
continuous on $(-1,1)$. Suppose $G=T_\lambda F$ almost everywhere,
$F,G$ are bounded, and $\int_{-1}^1w|G'|^2<\infty$.
Then $F'$ has a locally absolutely continuous representative, $wF'$ is
bounded, all integrals in Lemma~\ref{lem:mixed-boundary-square} are finite,
and that identity holds. No endpoint values of $F$ are required.
\end{lemma}

\begin{proof}
\uses{lem:ac-integration-by-parts}
\leanok
The equation gives $wF'=(x+i\lambda)F+iG$. On each compact subinterval
the right side divided by $w$ is continuous and absolutely continuous.
The fundamental theorem of calculus identifies it with the classical
derivative everywhere in the interior. Differentiating gives
$wF''-3xF'-i\lambda F'=iG'+F$, whose weighted square is integrable.
Also $T_0F=G+\lambda F$ is bounded. The difference of the integrands
is the derivative of
\[
 -w\left(x|F|^2+2\operatorname{Re}(iG\bar F)\right)
 = xw|F|^2-2w^2\operatorname{Re}(F'\bar F).
\]
The first expression is $O(w)$ at both ends by boundedness of $F,G$.
Integrate on compact subintervals and exhaust $(-1,1)$.
\end{proof}

\begin{lemma}[Stream trace on a mixed Neumann edge]
\label{lem:mixed-neumann-stream-trace}
\uses{def:mixed-quotient-data, def:h1-interval}
\lean{NoEmbeddedEigenvalues.mixed_neumann_stream_trace}
\leanok
Let $g=\psi|_N$, $x=\tanh t$, $w=1-x^2$, and $g=\sqrt w\,G$.
Then $g$ and $G$ are locally absolutely continuous, $g'\in L^1(dt)$,
$G$ is bounded, and $\int_{-1}^1w|G'|^2\,dx<\infty$.
\end{lemma}

\begin{proof}
\uses{lem:trace-h2-boundary, lem:arc-diagonal, lem:boundary-form-xi}
\leanok
The quotient $\phi=\psi/h$ is continuous on the closed cell and belongs to
$H^2$ in the cell, so its boundary restriction is the Sobolev trace
and belongs to $H^1(da)$ on the edge.
Along either directed Neumann edge $h=\sech t$ and $|dx|=h\,|da|$, so
\[
 G=\phi|_N,\qquad \int_{-1}^1w|G'|^2\,dx=\int_N h|\phi_a|^2\,da.
\]
Also $g'=-xhG+hwG'$. The first term is integrable in time since $G$ is
bounded and $h\in L^1(dt)$; the second is integrable since
$\int |hwG'|\,dt=\int_{-1}^1\sqrt w\,|G'|\,dx$.
\end{proof}

\begin{lemma}[Critical trace weight on an inward strip]
\label{lem:mixed-neumann-strip-weight}
\uses{def:sobolev-local}
\lean{NoEmbeddedEigenvalues.HasTransverseH1Trace.left_memLp_div_sqrt,
  NoEmbeddedEigenvalues.mixed_neumann_strip_memLp_flow,
  NoEmbeddedEigenvalues.mixed_neumann_inward_trace_memLp}
\leanok
On an inward strip $I\times(0,w)$, let $T$ be the canonical trace
of $F$ at normal coordinate $0$. Suppose $F/k$ and the weak normal
derivative of $F$ are square integrable, for a measurable tangential
weight $k>0$ such that $\delta k\le w$ for some $\delta>0$.
Then $T/\sqrt{k}\in L^2(I)$. On a full Neumann strip, $k=\cos a$
therefore gives $T\in L^2(dt)$ in flow time. In particular this applies to
the mixed vorticity on a Sobolev cell collar with the bulk quotient weight.
\end{lemma}

\begin{proof}
\uses{lem:trace-l2, lem:h1-interval-endpoints, lem:arc-diagonal,
  lem:edge-inward-pullback, lem:one-sided-trace-comparison}
\leanok
On each normal section, apply the endpoint trace inequality on the
interval of width $\delta k$. Divide by $k$ and enlarge the integrals
to the whole section to obtain
$|T|^2/k\le(2/\delta)\int|F/k|^2+2\delta\int|\partial_nF|^2$.
Integrate over the tangential variable. On the mixed arcs $|da|=\cos a\,dt$.
For the physical trace, pull back the smaller Sobolev collar to a full
inward strip, and compare to each cell-side window on their common
inward subwindow.
\end{proof}

\begin{lemma}[Physical vorticity trace on a mixed Neumann edge]
\label{lem:mixed-neumann-vorticity-trace}
\uses{def:mixed-quotient-data, def:sobolev-local}
\lean{NoEmbeddedEigenvalues.mixed_neumann_vorticity_trace}
\leanok
For a mixed eigenfunction with $\lambda\ne0$, the one-sided vorticity trace
on a Neumann edge has a representative $q$ in flow time with $q\in L^2(dt)$,
locally absolutely continuous and satisfying $q'=i\lambda q-g'$ almost
everywhere, where $g=\psi|_N$.
\end{lemma}

\begin{proof}
\uses{lem:mixed-cell-separatrix-h1, lem:mixed-neumann-stream-trace, lem:mixed-neumann-strip-weight,
  lem:edge-window-transport, lem:one-sided-trace-comparison, lem:transport-pair-eigen,
  lem:h1-interval-endpoints, lem:arc-diagonal}
\leanok
Choose a full inward strip in the cellwise $H^1$ collar. The critical
strip estimate gives a canonical trace $S(a)$ in $L^2(dt)$.
On each compact regular window the weak tangential equation gives
an $H^1$ representative of $S+g_0$, where $g_0$ is the continuous
stream trace in the side coordinate $a$. The derivative is
$i\lambda\kappa S/\cos a$, with $\kappa$ the sign of the flow direction.
The representatives agree pointwise on overlaps by continuity and
agreement almost everywhere; rational windows suffice to combine the
exceptional null sets. Compose the resulting representative with
$a(t)$, using $a'(t)=\kappa\cos a(t)$, and subtract $g(t)$.
This gives $q'=i\lambda q-g'$ and preserves both the $L^2(dt)$ class
and the canonical cell-side traces. No trace of the normal derivative
of $\omega$ is needed.
\end{proof}

\begin{lemma}[Decay of a square-integrable vorticity profile]
\label{lem:mixed-vorticity-profile-decay}
\uses{def:canonical-profile, def:h1-interval}
\lean{NoEmbeddedEigenvalues.mixed_vorticity_profile_decay}
\leanok
If $g'\in L^1(dt)$ and a locally absolutely continuous $q\in L^2(dt)$
satisfies $q'=i\lambda q-g'$, then $q$ tends to zero at both ends.
\end{lemma}

\begin{proof}
\uses{lem:ac-integration-by-parts}
The derivative of $e^{-i\lambda t}q(t)$ is $-e^{-i\lambda t}g'(t)$.
The integrating factor therefore gives a limit at each end. A nonzero
limit would keep $|q|$ bounded below by a positive constant on a half-line,
contradicting $q\in L^2(dt)$.
\leanok
\end{proof}

\begin{lemma}[Transformed canonical transport on a mixed Neumann edge]
\label{lem:mixed-neumann-canonical-transport}
\uses{def:canonical-profile, def:h1-interval}
\lean{NoEmbeddedEigenvalues.mixed_neumann_canonical_transport}
\leanok
Write $g=\sech t\,G(\tanh t)$ for the stream trace associated to a
quotient on a Neumann edge. Let $\lambda\ne0$ and suppose $G$ is bounded
on $(-1,1)$, $g$ is locally absolutely continuous, $g'\in L^1(dt)$, and the
Melnikov integral $\int e^{-i\lambda t}g'(t)\,dt$ vanishes. For the canonical
unknown $z$, the function $F$ defined by $z=\sech t\,F(\tanh t)$ is locally
absolutely continuous and bounded on $(-1,1)$, and $G=T_\lambda F$ a.e.
\end{lemma}

\begin{proof}
\uses{lem:canonical-profile-ode, lem:canonical-profile-ac, lem:canonical-profile-decay,
  lem:boundary-form-xi}
Local absolute continuity follows from that of $z=-({\Omega}+g)/\lambda$
by smooth rescaling and reparametrization on each compact interval.
Since $g=O(e^{-|t|})$, the two-sided canonical profile bound gives
$\Omega=O(e^{-|t|})$ and hence $F=(\cosh t)z=O(1)$.
Write $z'=i(\lambda z+g)$ in $x=\tanh t$ using $dx/dt=w$ to obtain
$G=T_\lambda F$.
\leanok
\end{proof}

\begin{lemma}[Boundary square for eigenfunction traces]
\label{lem:mixed-boundary-square-traces}
\uses{def:mixed-quotient-data, def:canonical-profile,
  lem:mixed-boundary-square-weighted, def:h1-interval}
\lean{NoEmbeddedEigenvalues.mixed_boundary_square_traces}
\leanok
Let $\omega$ be a mixed-sector eigenfunction with $\lambda\ne0$.
On a Neumann edge, use directed flow time $t$, put $x=\tanh t$ and
$w=1-x^2$, and write the canonical transport pair as
$z=\sqrt wF$, $\psi|_N=\sqrt wG$.
Then $F,G$ are locally absolutely continuous, $G$ is bounded,
$\int w|G'|^2<\infty$, $F$ and $wF'$ are bounded, and the identity of
Lemma~\ref{lem:mixed-boundary-square} holds. The physical vorticity
trace on the edge agrees almost everywhere with the canonical profile.
\end{lemma}

\begin{proof}
\uses{lem:mixed-neumann-stream-trace, lem:mixed-neumann-vorticity-trace,
  lem:mixed-vorticity-profile-decay, lem:canonical-profile-unique, lem:canonical-profile-ode,
  lem:mixed-neumann-canonical-transport, lem:mixed-boundary-square-weighted}
The decay at $+\infty$ identifies the physical trace with the canonical
profile by uniqueness. Decay at $-\infty$ then gives the Melnikov
condition. This supplies the bounded transformed unknown and the
first-order equation; apply the weighted boundary-square identity
with the Neumann stream energy.
\leanok
\end{proof}

\begin{remark}[Digest note: no elliptic-center input]
The manuscript invokes global $H^1$ regularity here. The critical
trace uses the inward collar from Lemma~\ref{lem:mixed-cell-separatrix-h1}
and the bulk quotient weight from Definition~\ref{def:mixed-quotient-data}.
Neither across-separatrix gluing nor elliptic-center regularity enters
the trace argument. Proposition~\ref{prop:mixed-quotient-domain}
is still needed to construct the supplied quotient data.
\end{remark}

\begin{definition}[Trace mass and static form]
\label{def:mixed-K}
\uses{def:mixed-field, prop:mixed-quotient-domain}
\lean{NoEmbeddedEigenvalues.mixedStaticForm,
 NoEmbeddedEigenvalues.mixedTraceMassCutoff,
 NoEmbeddedEigenvalues.mixedTraceMassDensity,
 NoEmbeddedEigenvalues.mixedVirialDensity}
\leanok
With $\psi=h\phi$, define
\[
  M_{\mathrm{tr}}[\phi]=2\operatorname{Re}\int_\Cell
     \frac{u\phi_a-v\phi_b}{p}\bar\phi\,da\,db,\qquad
  K[\phi]=J_W[h\phi]-M_{\mathrm{tr}}[\phi]-\sum_{P\text{ corner}}|\phi(P)|^2,
\]
where $J_W[\psi]=\int_\Cell(f|\omega|^2+\ip{M\nabla\psi}{\nabla\psi})\,dx\,dy$.
The volume definition avoids taking a limit before its existence is known.
\end{definition}

\begin{lemma}[Signed trace-mass field]
\label{lem:mixed-trace-field}
\uses{def:mixed-field}
\lean{NoEmbeddedEigenvalues.mixed_trace_mass_field, NoEmbeddedEigenvalues.mixed_trace_mass_edge_flux}
\leanok
The field $V=\frac up\partial_a-\frac vp\partial_b$ has zero divergence in
the diagonal rectangle and satisfies $Vp=h-k$. Its outer normal flux is
$1/k$ on the regular Dirichlet sides and $-1/h$ on the regular Neumann sides.
\end{lemma}

\begin{proof}
\leanok
Differentiate the two components and use $u^2+h^2=v^2+k^2=1$.
For the second identity, $Vp=(v^2-u^2)/p=(h^2-k^2)/p$.
Evaluate the normal components on each side using the sign of its
outward normal.
\end{proof}

\begin{lemma}[Trace mass as a volume integral]
\label{lem:mixed-trace-mass}
\uses{def:mixed-K}
\lean{NoEmbeddedEigenvalues.mixed_trace_mass}
\leanok
The field $V=\frac up\partial_a-\frac vp\partial_b$ is divergence free in the
$(a,b)$ rectangle, with outer normal flux $1/k$ on $D$ and $-1/h$ on $N$.
The volume integral defining $M_{\mathrm{tr}}[\phi]$ is absolutely convergent
and equals the paired limit
\[
  \lim_{T\to\infty}\Big(\sum_D\int_{-T}^T|\phi|^2\,dt
                         -\sum_N\int_{-T}^T|\phi|^2\,dt\Big).
\]
\end{lemma}

\begin{proof}
\uses{prop:mixed-quotient-domain, lem:mixed-trace-field, lem:arc-diagonal}
Use the divergence theorem for $|\phi|^2 V$ on $p>\delta$.
On each original side the condition $p>\delta$ is exactly the flow-time
cutoff with $\delta=\operatorname{sech}T$. On $p=\delta$, the normal flux
has magnitude $|h-k|/|\nabla_0p|=O(\delta)$: the new arcs are near the
corners where $|\nabla_0p|$ is bounded below, and their lengths are
$O(\delta)$. This avoids a separate cancellation argument on circular
arcs. To pass to the limit in the volume integral, use boundedness of
$\phi$ on the closed cell, $\nabla_0\phi\in L^4$ by the rectangle Sobolev
embedding, and $1/p\in L^{4/3}$ near the corners.
\end{proof}

\begin{lemma}[Renormalized Neumann flux]
\label{lem:mixed-neumann-renormalization}
\uses{def:mixed-quotient-data, def:canonical-profile, lem:mixed-boundary-square-traces}
\lean{NoEmbeddedEigenvalues.mixedNeumannFluxCutoff,
  NoEmbeddedEigenvalues.mixed_neumann_flux_renormalized}
\leanok
For a mixed eigenfunction and either Neumann side $N$, set
$B_N(T)=\int_{-T}^T (|\psi_t|^2_{\mathrm{tan}}+
\operatorname{Re}(\Omega\bar\psi))\,dt$, where
$|\psi_t|_{\mathrm{tan}}$ is the magnitude of the physical tangential
derivative, i.e. $|g'(t)|/(\sqrt2\sech t)$ for $g=\psi|_N$. Then
\[
 \lim_{T\to\infty}\left(2 B_N(T)-\int_{-T}^T|\phi|^2\,dt\right)
 = \int_{-1}^1 w|G'|^2\,dx-2\operatorname{Re}\int_{-1}^1 T_0F\bar G\,dx
   -\sum_{P\in\partial N}|\phi(P)|^2 .
\]
\end{lemma}

\begin{proof}
\uses{lem:mixed-boundary-square-traces, lem:arc-diagonal, lem:canonical-profile-ode,
  lem:boundary-form-xi}
\leanok
Substitute $x=\tanh t$ and $g=\sqrt w\,G$ on the centered truncated
interval. The renormalized gradient density is
$|wG'-xG|^2/w-|G|^2/w=w|G'|^2-(x|G|^2)'$.
Integrating the total derivative gives the two endpoint contributions,
which tend to the corner values by continuity. The pairing transforms
using $\Omega=-\sqrt w\,T_0F$; the remaining integrands are integrable by
Lemma~\ref{lem:mixed-boundary-square-traces}.
\end{proof}

\begin{lemma}[Static identity]
\label{lem:mixed-K-identity}
\uses{def:mixed-K, lem:mixed-boundary-square-traces}
\lean{NoEmbeddedEigenvalues.mixed_static_identity}
\leanok
Every mixed eigenfunction satisfies
$0=K[\phi]+\sum_N\int_{-1}^1w|wF''-3xF'-i\lambda F'|^2$. In particular
$K[\phi]\le0$ (manuscript \texttt{(eq:mixed-K-identity)}).
\end{lemma}

\begin{proof}
\uses{prop:mixed-virial, lem:mixed-trace-mass, lem:mixed-neumann-renormalization,
  lem:mixed-boundary-square-traces}
\leanok
At each cutoff, twice the paired flux is the trace-mass cutoff minus
the sum of the two renormalized Neumann contributions. Compare its two
limits using Proposition~\ref{prop:mixed-virial} and
Lemmas~\ref{lem:mixed-trace-mass} and~\ref{lem:mixed-neumann-renormalization}.
The two corner sums account for all four corners. The boundary-square
identity then gives $0=K[\phi]$ plus the nonnegative sum of squares.
\end{proof}

\section{The second-order form}

\begin{lemma}[Rellich form of $K$]
\label{lem:mixed-rellich}
\uses{def:mixed-K}
\lean{NoEmbeddedEigenvalues.mixed_rellich}
\leanok
$K[\phi]=\int_\Cell\big(f|\omega|^2+|\nabla\psi|^2-2\operatorname{Re}((W\psi)\bar\omega)\big)\,dx\,dy
-\sum_N\int h|\phi_a|^2\,da$ (manuscript \texttt{(eq:mixed-static-rellich)}).
\end{lemma}

\begin{proof}
\uses{lem:mixed-field-algebra, cor:mixed-saddle-jet}
Apply the Rellich identity
$\Div\{2\operatorname{Re}((W\psi)\nabla\bar\psi)-W|\nabla\psi|^2\}
=(2\sym DW-\Div W\,I)\nabla\psi\cdot\nabla\bar\psi+2\operatorname{Re}((W\psi)\bar\omega)$.
The Dirichlet-edge current cancels its part of $M_{\mathrm{tr}}$; on a
Neumann edge
$\int(-|\psi_a|^2/h+|\psi|^2/h^3)\,da=\sum_{\text{endpoints}}|\phi|^2-\int h|\phi_a|^2\,da$
produces the corner values. The corner flux of the leading jet integrates to
zero by $\theta\mapsto\pi/2-\theta$.
\end{proof}

\begin{lemma}[Weighted interior form]
\label{lem:mixed-weighted-form}
\uses{lem:mixed-rellich}
\lean{NoEmbeddedEigenvalues.mixed_weighted_form}
\leanok
With $L_\mu=\partial_a^2-2\tan a\,\partial_a+\partial_b^2$ and
$C_0=-f+\frac{2u^2}{hp}=1-\Div_\mu W$,
\[
  K[\phi]=\int\mu\Big[\frac f2|L_\mu\phi|^2+2\sym DW\nabla_0\phi\cdot\nabla_0\bar\phi
  +(2W-\nabla_0C_0)\cdot\operatorname{Re}(\phi\nabla_0\bar\phi)
  +\Big(1-C_0-\frac f2\Big)|\phi|^2\Big]\,da\,db
\]
(manuscript \texttt{(eq:weighted-static-K)}).
\end{lemma}

\begin{proof}
Use $\Delta_0(h\phi)=h(L_\mu-1)\phi$ and
$\int|\nabla_0(h\phi)|^2=\int\mu(|\nabla_0\phi|^2+|\phi|^2)$, and expand.
\end{proof}

\begin{definition}[Second-order coefficients]
\label{def:mixed-coefficients}
\uses{def:cells, lem:cell-square, lem:diagonal-calculus}
\lean{NoEmbeddedEigenvalues.mixedOpenCell, NoEmbeddedEigenvalues.mixedQuarterCell,
 NoEmbeddedEigenvalues.mixedD, NoEmbeddedEigenvalues.mixedC, NoEmbeddedEigenvalues.mixedV,
 NoEmbeddedEigenvalues.mixedComparisonDiagonal,
 NoEmbeddedEigenvalues.mixedComparisonOffDiagonal,
 NoEmbeddedEigenvalues.mixedComparisonPotential, NoEmbeddedEigenvalues.mixedComparisonOperator}
\leanok
In the local notation of this chapter let
\[
  G=\begin{pmatrix}g&j\\j&g\end{pmatrix},\quad g=h^2+\frac{2hu^2}p,\quad
  j=\frac{2h^2uv}{p^2},\quad V=-\frac{h^2R}{p^2}+\frac{2hu^2}p-\frac{2u^4}R,
\]
with $u^4/R$ assigned its continuous value $0$ at the center, and let
$\mathcal P=-\Div(G\nabla)+V$ in the $(a,b)$ variables.
\end{definition}

\begin{lemma}[Bounded second-order coefficients]
\label{lem:mixed-coefficient-bounds}
\uses{def:mixed-coefficients}
\lean{NoEmbeddedEigenvalues.mixed_comparison_coefficient_bounds}
\leanok
On the open diagonal-coordinate cell, $0\le-V\le4$, $g\le3$ and $|j|\le2$.
\end{lemma}

\begin{proof}
\leanok
Write $-V=(h\sqrt R/p-u^2/\sqrt R)^2+(u^2/\sqrt R)^2$ for $R>0$;
at $R=0$ the potential is zero. For the upper bound drop the negative
cross term and use $h\le p$ and $0\le u^2\le R\le2$. The bounds on $g$
and $j$ follow from $h^2\le1$, $h/p\le1$ and $|uv|\le1$.
\end{proof}

\begin{lemma}[Completion of the square]
\label{lem:mixed-completion}
\uses{lem:mixed-weighted-form, def:mixed-coefficients, def:sobolev-local}
\lean{NoEmbeddedEigenvalues.mixedSecondOrderForm, NoEmbeddedEigenvalues.mixed_completion,
 NoEmbeddedEigenvalues.mixed_second_order_form_bounded}
\leanok
\[
  K[\phi]=\int_\Cell\frac{\mu f}2\Big|L_\mu\phi+\frac{C_0}f\phi\Big|^2\,da\,db+Q_2[\phi],
\]
\[
  Q_2[\phi]=\int_\Cell\big(G\nabla_0\phi\cdot\nabla_0\bar\phi+V|\phi|^2\big)\,da\,db
  +\sum_N\int h|\phi|^2\,da .
\]
All terms are finite, and $Q_2$ is continuous on $H^1(\Cell)$ (manuscript
\texttt{(eq:static-second-completion)}).
\end{lemma}

\begin{proof}
\uses{prop:mixed-quotient-domain, lem:hardy-1d, lem:mixed-coefficient-bounds}
Integrate the drift by parts, complete the square, and use
$2\int\mu W\cdot\operatorname{Re}(\phi\nabla_0\bar\phi)
=-\int\mu\Div_\mu W|\phi|^2+\sum_N\int h|\phi|^2\,da$. Finiteness:
$\mu f$ and $\mu C_0^2/f=\mu f-4hu^2/p+4u^4/R$ are bounded, and
$L_\mu\phi=2\omega/h+\phi$. The corner currents vanish by angular Poincar\'e
applied to $\nabla_0\phi$.
\end{proof}

\begin{lemma}[Positivity of the second-order matrix]
\label{lem:mixed-G-positive}
\uses{def:mixed-coefficients}
\lean{NoEmbeddedEigenvalues.mixed_comparison_matrix_positive}
\leanok
$G$ is positive definite in the open cell:
\[
  \det G=\frac{h^2}{p^4}\big[h^6+4h(2-h^2)k+(4+8h^2-6h^4)k^2+4hk^3+h^2k^4\big]>0 .
\]
\end{lemma}

\begin{proof}
\leanok
Expand $g^2-j^2$ using $u^2=1-h^2$ and $v^2=1-k^2$.
The diagonal entry $g$ and every coefficient in the bracket are positive
for $0<h,k\le1$.
\end{proof}

\section{Positive comparisons}

\begin{definition}[Comparison form on the quarter cell]
\label{def:comparison-form}
\uses{def:mixed-coefficients, def:sobolev-local}
\lean{NoEmbeddedEigenvalues.mixedQuarterForm,
 NoEmbeddedEigenvalues.MixedQuarterFormDomain,
 NoEmbeddedEigenvalues.MixedQuarterZeroTrace}
\leanok
On $Q=(0,\pi/2)^2$, for $z\in H^1(Q)$, let
\[
  \mathcal Q[z]=\int_Q\big(G\nabla z\cdot\nabla\bar z+V|z|^2\big)\,da\,db
  +\int_0^{\pi/2}h|z(a,\pi/2)|^2\,da .
\]
(The manuscript writes $\mathsf G=G$ and $U=-V$ here.)
\end{definition}

\begin{definition}[Robin multiplier with logarithmic derivative bound]
\label{def:mixed-robin-multiplier}
\lean{NoEmbeddedEigenvalues.MixedRobinMultiplier}
\leanok
A Robin multiplier $\rho(h,k)$ is continuous on $[0,1]^2$, with
$0<m\le\rho\le M$ there, and smooth on an open neighbourhood of
$(0,1]\times[0,1]$. On this strip it satisfies
$|\rho_h|+|\rho_k|\le C(1+|\log h|)$ for some $C\ge0$.
\end{definition}

\begin{lemma}[Pointwise mixed Picone inequality]
\label{lem:mixed-robin-picone}
\uses{def:mixed-robin-multiplier, def:mixed-coefficients}
\lean{NoEmbeddedEigenvalues.mixed_robin_picone}
\leanok
Let $\rho$ be a Robin multiplier and take $g=u\rho$ or $g=v\rho$.
If $\mathcal Pg\ge ch^2g$ in $Q$, then for $z\in C^1(Q)$,
\[
 ch^2|z|^2+\operatorname{div}J
 \le G\nabla z\cdot\nabla\bar z+V|z|^2,\qquad
 J=|z|^2G\nabla g/g .
\]
\end{lemma}

\begin{proof}
\uses{lem:picone, lem:mixed-G-positive}
\leanok
The multiplier is positive and smooth in the open quarter. Use the Picone
divergence identity and drop its nonnegative gradient square.
\end{proof}

\begin{lemma}[Regular conormal currents]
\label{lem:mixed-robin-currents}
\uses{def:mixed-robin-multiplier, def:mixed-coefficients}
\lean{NoEmbeddedEigenvalues.mixed_robin_ground_logPartials,
 NoEmbeddedEigenvalues.mixed_robin_top_current,
 NoEmbeddedEigenvalues.mixed_robin_even_axis_current}
\leanok
The components of $\nabla g/g$ are
$X=\mathbf1_{g=u\rho}\,h/u-u\rho_h/\rho$ and
$Y=\mathbf1_{g=v\rho}\,k/v-v\rho_k/\rho$ on the regular strip in $(a,b)$,
away from the odd axis. On $b=\pi/2$, $0<a<\pi/2$, the sum of
the top current and the Robin density is
\[
 J_b+h|z|^2
 =\big((1+2\mathbf1_{g=u\rho})h\rho
 -2(1-h^2)\rho_h-(2-h^2)\rho_k\big)|z|^2/\rho .
\]
The conormal current on the regular even axis is zero.
\end{lemma}

\begin{proof}
\leanok
Differentiate the sine and cosine factors by the chain rule. On the top,
the diagonal entry is $2-h^2$ and the off-diagonal entry is $2u$.
On the even axis the off-diagonal entry and the even normal derivative
both vanish.
\end{proof}

\begin{lemma}[Mixed comparison on a regular rectangle]
\label{lem:mixed-robin-rectangle}
\uses{def:mixed-coefficients, def:mixed-robin-multiplier}
\lean{NoEmbeddedEigenvalues.mixed_robin_rectangle}
\leanok
Suppose $\mathcal Pg\ge ch^2g$ in $Q$ and the Robin numerator is
nonnegative on the top. Let $R=(a_0,A)\times(b_0,\pi/2)$, with
$0\le a_0<A<\pi/2$, $0\le b_0<\pi/2$, and with the closed rectangle
disjoint from the odd axis. For a globally $C^2$ test function $z$,
\[
 c\int_R h^2|z|^2 + \int_{b_0}^{\pi/2}(J_a(A,b)-J_a(a_0,b))\,db
       -\int_{a_0}^A J_b(a,b_0)\,da
 \le
 \int_R (G\nabla z\cdot\nabla\bar z+V|z|^2)
       +\int_{a_0}^A h|z(a,\pi/2)|^2\,da .
\]
\end{lemma}

\begin{proof}
\uses{lem:mixed-robin-picone, lem:mixed-robin-currents}
\leanok
Integrate the pointwise inequality. Fubini and the fundamental theorem
of calculus give the four signed normal currents. Combine the top
current with the Robin term and discard their nonnegative sum.
\end{proof}

\begin{lemma}[Comparison lemma with logarithmic weights]
\label{lem:robin-comparison}
\uses{def:comparison-form, def:mixed-robin-multiplier, def:sobolev-local}
\lean{NoEmbeddedEigenvalues.mixed_robin_comparison, NoEmbeddedEigenvalues.mixedRobinNumerator}
\leanok
Let $\rho$ be a Robin multiplier. The Robin inequalities below are only
on the regular edge $k=0$, $0<h\le1$. If $g=v\rho$
satisfies $\mathcal Pg\ge ch^2g$ in $Q$ and
$h\rho-2(1-h^2)\rho_h-(2-h^2)\rho_k\ge0$ on $k=0$, then
$\mathcal Q[z]\ge c\int_Qh^2|z|^2$ for all $z\in H^1(Q)$ with zero Sobolev
trace on $b=0$. The same holds for zero Sobolev trace on $a=0$ if $g=u\rho$
satisfies $\mathcal Pg\ge ch^2g$ and
$3h\rho-2(1-h^2)\rho_h-(2-h^2)\rho_k\ge0$ on $k=0$
(manuscript \texttt{(lem:logarithmic-robin-comparison)}).
\end{lemma}

\begin{proof}
\uses{lem:mixed-robin-rectangle, lem:mixed-robin-currents, lem:trace-l2}
Apply the regular rectangle inequality with $h>\epsilon$ and at distance
$\delta$ from the odd axis. For smooth tests vanishing on that axis, its
cutoff current is $O(\delta)$ for fixed $\epsilon$. The current on
$h=\epsilon$ is $O(\epsilon(1+|\log\epsilon|)^2)$ after integration along
the side. Extend the smooth-test inequality by density with the zero
Sobolev trace (odd reflection and mollification).
\end{proof}

\begin{definition}[Odd-in-$b$ comparison function]
\label{def:odd-b-comparison}
\uses{def:mixed-coefficients}
\lean{NoEmbeddedEigenvalues.mixedOddBMultiplier, NoEmbeddedEigenvalues.mixedOddBComparison}
\leanok
Let $w_v=vF_v(h,k)$ with
\[
  F_v=1+\frac{4132k^2+2167k^3-2011hk-7640hk^2-226h^2+3795h^2k-834h^3}{10000}
  +\frac{808hk+2505hk^2}{10000}\log h .
\]
\end{definition}

\begin{lemma}[Odd-in-$b$ multiplier and edge bounds]
\label{lem:mixed-v-multiplier-bounds}
\uses{def:odd-b-comparison}
\lean{NoEmbeddedEigenvalues.mixed_odd_b_multiplier_robin,
 NoEmbeddedEigenvalues.mixed_odd_b_multiplier_closed_bounds,
 NoEmbeddedEigenvalues.mixed_odd_b_multiplier_derivative_bound,
 NoEmbeddedEigenvalues.mixedOddBMultiplier_continuous,
 NoEmbeddedEigenvalues.mixedOddBMultiplier_contDiffAt}
\leanok
$F_v$ is continuous on the closed square and smooth wherever $h>0$.
There are strict bounds $\tfrac34<F_v<\tfrac53$ on $(0,1]\times[0,1]$,
and the non-strict bounds hold on the closed square. On the regular strip,
$|\partial_hF_v|+|\partial_kF_v|\le C(1+|\log h|)$ for some $C\ge0$.
On $k=0$, $0<h\le1$,
\[
  hF_v-2(1-h^2)\partial_hF_v-(2-h^2)\partial_kF_v\ge\tfrac{1789}{2500}h .
\]
\end{lemma}

\begin{proof}
\leanok
The quadratic form $2502h^2-7590hk+7640k^2$ is nonnegative, so the
polynomial part $F_0$ decreases in $h$. At $h=1$,
$10000(F_0(1,k)-0.894)=k(1784-3508k+2167k^2)\ge0$ by completing the
square. At $h=0$ the polynomial part is at most $1.6299$. The logarithmic
term is nonpositive, and $-h\log h\le e^{-1}<3/8$ gives its lower bound.
Write this term as a polynomial in $k$ times $h\log h$. The latter
extends continuously by zero and has derivative $1+\log h$ for $h>0$.
The bounds at $h=0$ follow from the polynomial part.
On the Robin side, differentiate at $k=0$ and divide the numerator by
$h>0$. The log term is nonnegative, and $h,h^2,h^3\le1$ give the margin.
\end{proof}

\begin{lemma}[Odd-in-$b$ comparison]
\label{lem:mixed-v-comparison}
\uses{def:odd-b-comparison}
\lean{NoEmbeddedEigenvalues.mixed_odd_b_interior}
\leanok
$\mathcal Pw_v\ge\tfrac1{100}h^2w_v$ in $Q$
(manuscript \texttt{(lem:mixed-v-comparison)}).
\end{lemma}

\begin{proof}
\uses{lem:mixed-v-residual-identity, cert:odd-b}
\leanok
Use the interior residual identity and the nonnegativity of $v\mathcal R_v$.
\end{proof}

\begin{definition}[Odd-in-$a$ comparison function]
\label{def:odd-a-comparison}
\uses{def:mixed-coefficients}
\lean{NoEmbeddedEigenvalues.mixedOddAExponent, NoEmbeddedEigenvalues.mixedOddAComparison}
\leanok
Let $w_u=ue^{F_u}$, where, with $S=h+k$,
\begin{align*}
1000F_u={}&721k^2-107k^3+60k^4+85k^6
 +h(-519k-64k^2-26k^3-400k^5)+h^2(-98-79k^2+890k^4)\\
 &+h^3(781-583k-842k^3)+h^4(-240+677k+124k^2)+h^5(-263+37k)+53h^6\\
 &-\frac h2-4k-\frac{4h^3}{S^2}-\frac{2744h^6}{S^4}+\frac{1033h^8}{S^6}
  +\frac{224h^4}{S}+\frac{235h^8}{S^5}
\end{align*}
(manuscript \texttt{(lem:mixed-u-comparison)}).
\end{definition}

\begin{lemma}[Bounded exponent and first derivatives]
\label{lem:mixed-u-exponent-bounds}
\uses{def:odd-a-comparison}
\lean{NoEmbeddedEigenvalues.mixed_odd_a_exponent_bounds_regular,
 NoEmbeddedEigenvalues.mixedOddAExponent_continuousOn,
 NoEmbeddedEigenvalues.mixedOddAExponent_contDiffAt}
\leanok
$F_u$ is continuous on the closed square and smooth on $h+k>0$.
There is a finite constant bounding $|F_u|$, $|\partial_hF_u|$ and
$|\partial_kF_u|$ throughout $(0,1]\times[0,1]$.
\end{lemma}

\begin{proof}
\leanok
Each rational term is a constant times $h^m/(h+k)^n$ with $m>n$.
This is bounded by $h^{m-n}$ on the square, and hence extends by zero at
the corner. Both partial derivatives are linear combinations of such
quotients with numerator and denominator degrees satisfying $m'\ge n'$,
so they are bounded, including on the regular edge. This does not give a
continuous gradient at $(0,0)$.
\end{proof}

\begin{lemma}[Odd-in-$a$ comparison]
\label{lem:mixed-u-comparison}
\uses{def:odd-a-comparison}
\lean{NoEmbeddedEigenvalues.mixed_odd_a_robin,
 NoEmbeddedEigenvalues.mixed_odd_a_interior}
\leanok
$\mathcal Pw_u\ge\tfrac1{200}h^2w_u$ in $Q$, and on the regular edge
$k=0$, $0<h\le1$,
\[
  3h-2(1-h^2)\partial_hF_u-(2-h^2)\partial_kF_u\ge\tfrac h{250}
\]
(manuscript \texttt{(lem:mixed-u-comparison)}).
\end{lemma}

\begin{proof}
\uses{lem:mixed-u-residual-identity, lem:mixed-robin-identity, cert:odd-a, cert:robin}
\leanok
Multiply the nonnegative interior residual by the positive comparison
function, and apply the Robin sign to the regular-edge identity.
\end{proof}

\section{Conclusion}

\begin{lemma}[Parity splitting]
\label{lem:mixed-parity-split}
\uses{def:mixed-quotient-data, lem:mixed-completion, def:sobolev-local}
\lean{NoEmbeddedEigenvalues.mixed_parity_split, NoEmbeddedEigenvalues.mixed_parity_mass}
\leanok
Write $\phi=\phi_u+\phi_v$ with $\phi_u$ odd in $a$ and even in $b$, $\phi_v$
even in $a$ and odd in $b$ (possible since $\phi$ is centrally odd). Then
$Q_2[\phi]=Q_2[\phi_u]+Q_2[\phi_v]$.
The same splitting holds for the weighted mass $\int_\Cell h^2|\phi|^2\,da\,db$.
\end{lemma}

\begin{proof}
\uses{lem:axis-reflection-weak-jet, lem:affine-jet-pullback,
  lem:specified-jet-lp, lem:mixed-coefficient-bounds, lem:mixed-G-positive}
\leanok
The diagonal and scalar coefficients of $Q_2$ are even in both variables and
the off-diagonal coefficient is odd in each. Reflect the function and its
specified weak gradient in $a$. The sum of the two projected
densities is the average of the original and reflected densities by the
parallelogram identity. Use the same argument on each Robin edge and
for the weighted mass.
\end{proof}

\begin{lemma}[Admissibility of the explicit Robin multipliers]
\label{lem:mixed-robin-multipliers}
\uses{def:mixed-robin-multiplier, def:odd-a-comparison, def:odd-b-comparison}
\lean{NoEmbeddedEigenvalues.mixed_explicit_robin_multipliers}
\leanok
Both $e^{F_u}$ and $F_v$ are Robin multipliers.
\end{lemma}

\begin{proof}
\uses{lem:mixed-u-exponent-bounds, lem:mixed-v-multiplier-bounds}
\leanok
Continuity of $F_u$ on the compact square gives a positive lower bound
and a finite upper bound for its exponential. The chain rule gives
bounded derivatives on the regular strip, which implies the logarithmic
bound there. For $F_v$, apply the multiplier bounds directly.
\end{proof}

\begin{lemma}[Coercivity on the quarter cell]
\label{lem:mixed-quarter-coercivity}
\uses{def:comparison-form, def:sobolev-local}
\lean{NoEmbeddedEigenvalues.mixed_quarter_coercivity}
\leanok
For quarter-form data with zero Sobolev trace on either odd axis,
$\mathcal Q[z]\ge \frac1{200}\int_Q h^2|z|^2$.
\end{lemma}

\begin{proof}
\uses{lem:robin-comparison, lem:mixed-robin-multipliers,
  lem:mixed-u-comparison, lem:mixed-v-comparison, lem:mixed-v-multiplier-bounds}
\leanok
Apply the Robin comparison with the multiplier for the specified odd axis.
In the exponential case factor out the positive $e^{F_u}$ from the
Robin numerator; in the other case weaken $1/100$ to $1/200$.
\end{proof}

\begin{lemma}[Quarter reduction with canonical traces]
\label{lem:mixed-fold-quarter}
\uses{def:comparison-form, lem:mixed-completion, def:sobolev-local}
\lean{NoEmbeddedEigenvalues.mixed_fold_to_quarter}
\leanok
An $H^1$ function continuous on the closed cell, odd in one axis and even
in the other, restricts to quarter-form data. Its Robin trace is its
continuous edge value and its trace vanishes on the odd axis. Its full
weighted mass is four times the quarter mass and its full second-order
form is four times the quarter form.
\end{lemma}

\begin{proof}
\uses{lem:axis-reflection-weak-jet, lem:canonical-transverse-trace,
  lem:mixed-coefficient-bounds, lem:mixed-G-positive}
Weak gradient parity gives reflection invariance of the form density;
the axes have measure zero. The trace agrees with the continuous
representative on the closed normal sections, and the Robin edge
integrals agree by reflection.
\end{proof}

\begin{lemma}[Second-order coercivity on the cell]
\label{lem:mixed-second-order-coercivity}
\uses{def:mixed-quotient-data, lem:mixed-completion}
\lean{NoEmbeddedEigenvalues.mixed_second_order_coercivity}
\leanok
$Q_2[\phi]\ge\frac1{200}\int_\Cell h^2|\phi|^2\,da\,db$ for the mixed quotient.
\end{lemma}

\begin{proof}
\uses{lem:mixed-parity-split, lem:mixed-fold-quarter, lem:mixed-quarter-coercivity}
\leanok
Split into the two parity components and apply the quarter bound to each.
The reflection factors are the same for the form and the mass.
\end{proof}

\begin{lemma}[Cartesian and diagonal stream mass]
\label{lem:mixed-stream-mass}
\uses{def:mixed-quotient-data, lem:cell-square}
\lean{NoEmbeddedEigenvalues.MixedQuotientData.stream_mass}
\leanok
$\int_\Cell|\psi|^2\,dx\,dy=2\int_\Cell h^2|\phi|^2\,da\,db$ for the mixed quotient.
\end{lemma}

\begin{proof}
\uses{lem:diagonal-l2-transport}
\leanok
Change coordinates and use the stream factorization almost everywhere.
\end{proof}

\begin{theorem}[Static coercivity]
\label{thm:mixed-coercivity}
\uses{lem:mixed-completion}
\lean{NoEmbeddedEigenvalues.mixed_static_coercivity}
\leanok
For the mixed quotient,
$K[\phi]\ge\frac1{200}\int_\Cell h^2|\phi|^2\,da\,db=\frac1{400}\int_\Cell|\psi|^2\,dx\,dy$
(manuscript \texttt{(eq:mixed-static-coercivity)}).
\end{theorem}

\begin{proof}
\uses{lem:mixed-completion, lem:mixed-second-order-coercivity, lem:mixed-stream-mass}
\leanok
Drop the nonnegative square from the completion and use the second-order
coercivity bound. The Cartesian mass introduces a factor of two.
\end{proof}

\begin{theorem}[Mixed exclusion]
\label{thm:mixed-exclusion}
\uses{def:eigen-problem, def:sectors}
\lean{NoEmbeddedEigenvalues.mixed_exclusion}
\leanok
If $\omega$ solves the eigenvalue problem with $\lambda\ne0$ and
$S_1^*\omega=-\omega$, $S_2^*\omega=\omega$, then $\omega=0$.
\end{theorem}

\begin{proof}
\uses{thm:mixed-coercivity, lem:mixed-K-identity, prop:mixed-quotient-domain,
  lem:stream-sector, lem:ee-cell-vanishing, lem:canonical-stream-poisson}
\leanok
$K[\phi]\le0$ by Lemma~\ref{lem:mixed-K-identity} and
$K[\phi]\ge\frac1{400}\int_\Cell|\psi|^2$ by Theorem~\ref{thm:mixed-coercivity}.
Thus the stream mass vanishes and $\psi=0$ almost everywhere on $\Cell$.
The stream inherits the vorticity's parities, so vanishing extends to the
torus by Lemma~\ref{lem:ee-cell-vanishing}. The Fourier Poisson relation
then gives $\omega=0$.
\end{proof}

\begin{corollary}[The other mixed sector]
\label{cor:mixed-other-exclusion}
\uses{def:eigen-problem, def:sectors}
\lean{NoEmbeddedEigenvalues.mixed_other_exclusion}
\leanok
If $\omega$ solves the eigenvalue problem with $\lambda\ne0$ and
$S_1^*\omega=\omega$, $S_2^*\omega=-\omega$, then $\omega=0$.
\end{corollary}

\begin{proof}
\uses{thm:mixed-exclusion, lem:sector-swap, lem:symmetries-commute}
\leanok
$R^*\omega$ solves the eigenvalue problem with the same $\lambda$
(Lemma~\ref{lem:symmetries-commute}) and lies in sector $(-,+)$
(Lemma~\ref{lem:sector-swap}), so $R^*\omega=0$ by
Theorem~\ref{thm:mixed-exclusion}. Since $R^*$ is an isometry, $\omega=0$.
\end{proof}
